\section{Results, scope, and the repository baseline}
\label{sec:overview}

The useful common framework is not a blanket principle saying that
computable mathematics belongs to \iPRA. It is a disciplined interface
between finite observations, numerical bounds, and proofs. Many of the
repository's reports already use parts of this interface. We consolidate
those parts into compositional theorems, strengthen the interface where
nested constructions need it, and use it to prove specific new extensions.

The inspected repository revision is
\begin{center}
\texttt{94000996e69ab0c527393a1d02a4a037a4eb0493}.
\end{center}
The directory inventory contains 48 delivered proof packets. Its own
README explicitly says that placement did not rerun or validate all the
packets, and that many Lean files are outlines \cite{repo-index}. Our
audit likewise does not certify every theorem in that collection. The
provenance ledger records the particular files used, their repository
blob hashes, and the source locations of the obstacles addressed below.

\subsection{What is established here}

\begin{longtable}{@{}>{\raggedright\arraybackslash}p{0.24\linewidth}>{\raggedright\arraybackslash}p{0.46\linewidth}>{\raggedright\arraybackslash}p{0.24\linewidth}@{}}
\toprule
Result & Mathematical content & Dependence and status\\
\midrule
\endhead
Finite observation compilation & Variable-depth operations are evaluated
by a bounded input table, a backward demand schedule, and a forward
numerical recursion. Finite contracts propagate by quantifier-free
induction. & General human proof; selected locality and compilation
lemmas checked in Lean.\\[5pt]
Uniform polynomial factorization & For each fixed PR coefficient evaluator,
one PR program factors monic polynomials uniformly in numerical degree,
returning coherent complex names and explicit coefficient-error bounds.
& Uses the existing finite rational descent lemma; all new normalization,
deflation, and iteration estimates proved here.\\[5pt]
Derived Bohr--Mollerup modulus & Positive rational multiplicative convexity
gives a local Lipschitz bound from two endpoint samples, so the candidate
modulus is no longer independent input. & Direct elementary proof;
canonical extension of the rational presentation.\\[5pt]
Finite-net extrema and coherent minimizers & A net and a supplied modulus
give an infimum name; quantitative uniqueness gives a coherent optimizer.
Strong convexity provides an explicit uniqueness modulus. & Direct
finite constructions and complete quantitative proofs.\\[5pt]
Superelliptic section bound & For $y^r=f(x)$, $\deg f=m>rs$, a proper
subspace meets $(1,x,\ldots,x^s,y)$ in at most $m$ distinct point images.
& Unconditional finite algebra. A tuple-local $M$-subspace cover then
gives the conditional cardinality bound $mM$.\\
\bottomrule
\end{longtable}

The new factorization theorem closes the fixed-degree restriction in the
existing FTA discussion, with its coefficient evaluator still fixed in the
correct first-order sense. The new convexity argument answers an explicit
further-research question in the Bohr--Mollerup report. The Diophantine
result discharges one geometric input for a defined family; it does not
supply the quantitative approximation theorem that must produce the
covers.

\subsection{The central organization}

The following separation makes mathematical reuse safe.
\begin{enumerate}
\item \textbf{Representation.} State exactly what finite observations
      are available. Exact algebraic data, Cauchy names, formal
      coefficients, and names with apartness certificates are distinct
      interfaces.
\item \textbf{Construction.} Give primitive recursive numerical
      procedures with an explicit demand on their input observations.
\item \textbf{Verification.} State which finitely many input contracts
      establish the requested output facts, and bound their demand
      separately from the construction's own queries.
\item \textbf{Arithmetic proof.} Fix the numerical output challenge,
      compile all its verification demands, and induct on a Boolean
      property of the resulting finite data.
\item \textbf{Interpretation.} Only after this finite argument quantify
      over the arbitrary output challenge and read the result as an
      identity, a continuous object, a root, or a finiteness assertion.
\end{enumerate}

This organization is useful even when it does not remove a hypothesis.
For example, it pinpoints the analytic estimate still needed for a
stronger prime-distribution theorem, instead of hiding that estimate
inside an interface called ``effective asymptotics.'' It also explains
why a finite cardinality theorem does not automatically enumerate all
solutions.

\subsection{How to read the report}

Sections~\ref{sec:foundations} and~\ref{sec:compiler} give the shared
arithmetic language and the compilation theorem.
Section~\ref{sec:fta-extension} is the most detailed variable-depth
application. Section~\ref{sec:quantitative-analysis} treats precision,
extrema, uniqueness, and convexity. Section~\ref{sec:finite-sections}
provides the elementary algebraic zero count and its conditional
Diophantine interface. The final sections explain exactly which reported
obstacles remain, how to complete the Lean development, and what was
actually checked in this archive.
